\section{Managed Accounts, Virtual Portfolios, and TRS}
\label{sec:sec09}
\label{sec:managed}

A managed account is not a primitive. It is the composition of four that already exist
(\S\ref{sec:ledger}, \S\ref{sec:sec04}): a wallet partition that holds positions and produces
PnL; a mandate unit $u_{\mathrm{MA}}$ issued by the manager and held by the client, conserved by
the issuance law; a deterministic reset contract that observes performance, crystallises one net
cash flow, and resets its baseline; and the three-home state model of \S\ref{sec:sec04}, which
homes the mandate's terms, the shared observables, and the per-client economic state. A total
return swap is the same composition with the wallet partition placed in a virtual ledger and the
reset contract's settlement crossing into the real one. Both reduce to a single conservative cash
move whose magnitude is a function of observed performance and a recorded baseline; the worked
example (\S\ref{sec:sec09-example}) carries one account through subscription, trade, NAV,
crystallisation against a high-water mark, and redemption, with $\sum_w w(u)=0$ at every step.

\subsection{The mandate is a unit}
\label{sec:sec09-mandate}

A managed account is opened by issuing the mandate unit $u_{\mathrm{MA}}$ from the manager $w_M$
to the client $w_C$, giving $w_M(u_{\mathrm{MA}})=-1$, $w_C(u_{\mathrm{MA}})=+1$,
$\sum_w w(u_{\mathrm{MA}})=0$. This is a real two-named-wallet issuance with a real issuer and
holder, not a self-reflexive sentinel; conservation holds by the issuance law (P1).

A \texttt{Move} is one conserved transfer of a unit between two named wallets:
$\text{from}\mathbin{-}\!=q$, $\text{to}\mathbin{+}\!=q$. Conservation (P1) is by construction,
not by check --- the per-wallet deltas are group negatives, so they sum to the identity. The
mandate issuance is the canonical instance.

\begin{lstlisting}[language=Haskell]
data Move = Move
  { mFrom :: WalletId, mTo :: WalletId, mUnit :: UnitId, mQty :: Qty }  -- q > 0

-- The per-wallet deltas a move induces; their sum is the conservation witness.
moveDelta :: Move -> [(WalletId, Qty)]
moveDelta (Move s d _ q) = [(s, qneg q), (d, q)]   -- qneg q <> q = mempty : conserved

-- Mandate issuance: a real two-named-wallet transfer of u_MA, manager -> client,
-- q = 1. A real issuer and holder, not a self-reflexive sentinel.
issueMandate :: WalletId -> WalletId -> UnitId -> Move
issueMandate manager client uMA = Move manager client uMA (Qty 1)
\end{lstlisting}

\begin{invariant}[Mandate is a singleton bilateral contract]
\label{inv:singleton}
While the mandate is live,
\[
\operatorname{support}\big(w\mapsto w(u_{\mathrm{MA}})\big)=\{(w_M,-1),(w_C,+1)\}.
\]
\end{invariant}

Conservation alone does not establish Invariant~\ref{inv:singleton}: P1 admits $(-5,+5)$ or a
three-way split. Identity is pinned by four existing primitives together: Invariant~\ref{inv:singleton}
for the active phase; the \texttt{Option} accessor distinguishing \texttt{None} (pre-issuance)
from \texttt{Some(zero)} (terminated, retained); a \textit{UnitStatus} lifecycle guard admitting
issuance only from an unissued status, so a terminated $u_{\mathrm{MA}}$ cannot be resurrected;
and C8 routing a re-mandate to a fresh $u'_{\mathrm{MA}}$ stamped \texttt{SetSupersededBy}.

\begin{principle}[Mandate is non-valued]
\label{prin:nonvalued}
$u_{\mathrm{MA}}$ is excluded from the domain of valuation: $P_t(u_{\mathrm{MA}})$ is undefined,
not zero. A price on $u_{\mathrm{MA}}$ would add $+1\cdot P_t(u_{\mathrm{MA}})$ to the client's
per-wallet NAV and double-count the underlying exposure; the term cancels only in the global sum.
\end{principle}

\paragraph{What is recorded where.}

\begin{longtable}{p{8.0cm}p{6.4cm}}
\toprule
\textbf{Datum} & \textbf{Home} \\
\midrule
Mandate text, fee schedule, benchmark identity, position limits, HWM/hurdle methodology,
crystallisation frequency & \textit{ProductTerms}$[u_{\mathrm{MA}}]$ (immutable, versioned,
append-only) \\
Lifecycle stage, \texttt{superseded\_by} & \textit{UnitStatus}$[u_{\mathrm{MA}}]$ (shared) \\
HWM value, entry NAV, accrued fees, mandate-breach flags, subscription/redemption cursor,
benchmark NAV at this wallet's inception & \textit{PositionState}$[w_C,u_{\mathrm{MA}}]$ \\
Current benchmark level (externally sourced, captured as a logged observation event) &
\textit{UnitStatus}$[u_{\mathrm{bench}}]$ \\
\bottomrule
\end{longtable}

The benchmark level is a shared cached observable, governed by the same discipline as every
\textit{UnitStatus} value (\S\ref{sec:sec04}): UnitStatus holds one value per unit, shared---read identically by every holder. Its value changes over time, but UnitStatus is not a separate
source of truth: the immutable event log is. Every change is caused by a logged event; the stored
value is overwritten in place only as a read cache. Replaying the events up to any point rebuilds
the exact value that held then, so nothing is lost by overwriting and there is no other way to
change the value. A value exists from registration onward.

The per-client economic scalars are homed in \textit{PositionState}, yet each is a fold of the
event log: whether they are stored or re-derived is escalation~E1 (\S\ref{sec:conclusion}).

\subsection{Subscription}
\label{sec:sec09-subscription}

Capital enters as a cash move from the client's external funding wallet $w_{\text{ext}}$, a
virtual wallet at the ledger boundary, into $w_C$, tagged \texttt{SUBSCRIPTION} and indexed by the
subscription/redemption cursor in \textit{PositionState}$[w_C,u_{\mathrm{MA}}]$. The tag is
load-bearing: it lets the net external flow over any interval be reconstructed as a projection
over the stream (\S\ref{sec:sec09-fees}). At subscription the entry NAV and the initial high-water
mark are set to the subscribed capital base; the high-water mark before subscription is
\texttt{None} (unrepresentable), never $0$ --- a $0$ mark would charge a performance fee on the
entire contributed capital at the first crystallisation. The benchmark NAV at inception is
recorded from \textit{UnitStatus}$[u_{\mathrm{bench}}]$ for performance reporting.

\subsection{Trading under mandate}
\label{sec:sec09-trading}

A trade is an ordinary two-legged transaction against a virtual market wallet $w_{\text{mkt}}$,
admitted only if it satisfies the mandate guards.

\begin{definition}[Mandate guard]
A quantitative mandate constraint is a precondition $g$ on the move-generation function. The
contract evaluates $g$ against the current state and the recorded price snapshot; if $g$ fails, no
moves are emitted and the state is unchanged.
\end{definition}

Guard determinism is price-relative. ``Same wallet state $+$ same trade $\Rightarrow$ same
decision'' is false for value-based limits --- leverage, concentration-by-value, currency caps ---
because holding positions fixed and doubling a price flips the same trade from admit to reject.
Determinism is recovered only when the guard is a pure function of
$(\textit{ProductTerms},\,\text{state},\,\text{trade},\,P_t)$ with $P_t$ an explicit,
snapshot-pinned input. Two totality obligations follow. A missing price ($P_t(u)=\texttt{None}$)
makes the guard fail closed, a typed reject, never a silent admit. A passive breach --- a limit
crossed by price appreciation with no admitted trade --- lies outside the domain of a move-time
guard and is detected only by a periodic valuation sweep. The pinned-snapshot requirement is
escalation~E2 (\S\ref{sec:conclusion}).

\subsection{Valuation and NAV}
\label{sec:sec09-nav}

The account NAV is the wallet valuation with the mandate excluded
(Principle~\ref{prin:nonvalued}):
\[
\mathrm{NAV}_t \;=\; \sum_{u\neq u_{\mathrm{MA}}} w_{C,t}(u)\,P_t(u).
\]
By state-sufficiency, $\mathrm{NAV}_t$ depends only on current balances, current unit state, and
current prices, not on history. Performance against the benchmark is computable from the inception
benchmark NAV and the current level in \textit{UnitStatus}$[u_{\mathrm{bench}}]$; it feeds
reporting, and feeds the fee only when the mandate's methodology names a benchmark hurdle.

\subsection{Fee logic}
\label{sec:sec09-fees}

\paragraph{Performance is net of capital flows.} Over a reset interval $[t_{k-1},t_k]$ the
performance is the gross NAV change less the net external flow:
\[
\mathrm{Perf}_k \;=\; \big(\mathrm{NAV}_{t_k}-\mathrm{NAV}_{t_{k-1}}\big)\;-\;
\mathrm{NetFlow}_{[t_{k-1},t_k]},\qquad
\mathrm{NetFlow}=\!\!\sum_{\substack{m\ \text{tagged sub/redemption}\\ t_{k-1}<m.t\le t_k}}\!\!\pm\,m.q.
\]
The gross change conflates performance with capital flows: by the PnL decomposition a mid-period
subscription is a flow, and crystallising it as performance would charge a fee on the client's own
contributed capital. $\mathrm{NetFlow}$ is a projection over the already-tagged cursor moves
(\S\ref{sec:sec09-subscription}); the correction adds no new coordinate.

\paragraph{Fee is not PnL settlement.} A single whole-$\mathrm{Perf}\to$beneficiary move is the
internal-desk case, where a book is swept to Treasury. A client managed account pays the manager
only a fee fraction; the gain stays in the client book. The fee handler is total over
$\operatorname{sign}(\mathrm{NAV}-\mathrm{HWM})$:
{\small\begin{align*}
\text{management fee } f^{m}_k &= r_{\text{m}}\cdot\Delta t_k\cdot \mathrm{NAV}^{\text{base}}_k
&&\text{\footnotesize(accrues on AUM regardless of sign),}\\
\text{performance fee } f^{p}_k &= r_{\text{p}}\cdot\max\!\big(0,\ \mathrm{NAV}^{\text{net}}_k-\mathrm{HWM}_{k-1}\big)
&&\text{\footnotesize(floored at zero; no clawback),}\\
\mathrm{HWM}_k &= \max\!\big(\mathrm{HWM}_{k-1},\ \mathrm{NAV}^{\text{net}}_k-f^{p}_k\big)
&&\text{\footnotesize(monotone ratchet; C11 writer \texttt{fee\_crystallise}),}
\end{align*}}
where $\mathrm{NAV}^{\text{net}}_k=\mathrm{NAV}_{t_k}-f^{m}_k$. On a loss period the performance
fee is $0$ and the high-water mark is unchanged, while the management fee still accrues.

\paragraph{Crystallisation is double-entry.} A fee is realised as a real cash move
$w_C\to w_M$, so it is conserved and carries its own contra-entry; a non-conserved stored scalar
would be single-entry and sit outside P1. The handler emits both legs atomically (C3): a
management-fee move and a performance-fee move, $w_C\to w_M$ in USD.

A signed amount $x$ crystallises into \emph{at most one} conserved cash move:
magnitude $|x|$, direction by $\operatorname{sign}(x)$, no move at $x=0$. Totality is the
\texttt{Maybe} --- the $x=0$ case is \texttt{Nothing}, never an illegal $q=0$ move. The fee
handler emits two such moves $w_C\to w_M$ atomically (C3) and ratchets the high-water mark
through the one canonical writer \texttt{WHwm :: FieldWrite 'FeeCrystallise} (reused verbatim
from the three-home model, \S\ref{sec:sec04}): the phantom index --- meaning the write names
its handler in its type --- makes a settle or trade handler bumping the HWM a compile error, so
the mark has exactly one author. \texttt{applyWrite (WHwm \_)} ratchets by \texttt{qmax}, so
$\mathrm{HWM}_k=\max(\mathrm{HWM}_{k-1},\mathrm{NAV}^{\text{net}}_k-f^p_k)$ holds by
construction, not by a runtime branch.

\begin{lstlisting}[language=Haskell]
-- Crystallise(magnitude, from, to): |x|, direction by sign(x), no move at x = 0.
crystallise :: Qty -> WalletId -> WalletId -> Maybe Move
crystallise (Qty x) from to
  | x > 0     = Just (Move from to usd (Qty x))            -- from -> to
  | x < 0     = Just (Move to   from usd (Qty (negate x))) -- direction flips
  | otherwise = Nothing                                    -- no move at zero

data FeeEvent = FeeEvent
  { fcMgmt :: Qty, fcPerf :: Qty, fcNavNet :: Qty }   -- f_m, f_p, NAV net of f_m

-- Two cash moves w_C -> w_M (mgmt, perf), atomic (C3), + the HWM ratchet write.
-- f_m, f_p >= 0; on a loss period f_p = 0, so crystallise returns Nothing and the
-- performance move is absent (no zero-quantity move).
crystalliseFees :: WalletId -> WalletId -> FeeEvent
                -> ([Move], FieldWrite 'FeeCrystallise)
crystalliseFees wC wM fc =
  ( [ m | Just m <- [ crystallise (fcMgmt fc) wC wM
                    , crystallise (fcPerf fc) wC wM ] ]
  , WHwm (fcNavNet fc <> qneg (fcPerf fc)) )   -- HWM_k = max(HWM_{k-1}, NAV^net - f_p)
\end{lstlisting}

\noindent At the worked example's Step~E,
\texttt{crystalliseFees wC wM (FeeEvent (Qty 5750) (Qty 28850) (Qty 1144250))} yields the
two legs \texttt{Move wC wM USD (Qty 5750)} and \texttt{Move wC wM USD (Qty 28850)} together
with \texttt{WHwm (Qty 1115400)}; applied to a row with $\mathrm{HWM}_{k-1}=1{,}000{,}000$ the
\texttt{qmax} ratchets the mark to $1{,}115{,}400$. On a loss period $f^p_k=0$ leaves one leg
and the mark unchanged.

\begin{theorem}[Fee zero-sum per crystallisation]
\label{thm:fee-zerosum}
$\sum_w \Delta_{\text{fee}}(\text{USD})=0$ at each crystallisation: the client pays exactly what
the manager receives, creating and destroying no value.
\end{theorem}
\noindent\textit{Proof.} Each leg is $w_C\mathbin{-}\!=q;\ w_M\mathbin{+}\!=q$; summing over
wallets gives $-q+q=0$ per leg, hence $0$ over the transaction. \qed

\paragraph{Quantity, baseline, residual.} Every settlement leg here is a \texttt{crystallise}
call (above): magnitude $|x|$, direction by $\operatorname{sign}(x)$, no move at $x=0$.
The reset baseline is the post-settlement capital base
$B_k=\mathrm{NAV}_{t_k}-(f^{m}_k+f^{p}_k)+\mathrm{NetFlow}$; reading the pre-settlement NAV would
claw the distributed performance back next period. The rounding residual
$\mathrm{Perf}_k-\operatorname{round}(\mathrm{Perf}_k)$ is retained in $w_C$ as a conserved
un-crystallised remainder, never dropped; it is the reconciling item between economic performance
and moved cash. Whether $B_k$, $\mathrm{HWM}_k$, and the accrued fee are stored or re-derived is
escalation~E1 (\S\ref{sec:conclusion}). Crystallising unrealised mark-to-market can drive the cash
wallet negative; conservation alone cannot distinguish a funded obligation from an insolvent
overdraft, a liveness gap recorded as escalation~E5 (\S\ref{sec:conclusion}).

\subsection{Segregation}
\label{sec:sec09-segregation}

\begin{theorem}[Segregation]
\label{thm:seg}
Client partitions are mutually isolated if and only if conservation is combined with locality and
capability scoping. A cross-client move is an illegal state ruled out by C4 under locality, not by
conservation.
\end{theorem}

Conservation alone is not segregation. A move $m=(w_A,w_B,\text{USD},q)$ between two different
clients' partitions is perfectly conservative ($\sum_w\Delta=0$), so P1 does not forbid it. Two
independent facts are at work: conservation (CONS, $\sum_w\Delta=0$) and locality (LOC,
$\operatorname{supp}(\Delta)\subseteq\{w_s,w_d\}$). Locality gives ``no effect beyond the named
wallets''; isolation of distinct clients additionally requires an authorisation predicate --- C4
capability scoping, forbidding a handler from naming or reading across a $(w,u_{\mathrm{MA}})$
boundary it has no capability for. Segregation is therefore a theorem under CONS $\wedge$ LOC
$\wedge$ C4, supplying the logical-segregation component of CASS~6 / MiFID~II~Art.~16(8). Legal
segregation --- distinct custodial accounts, trust or nominee structure --- is outside scope: the
ledger evidences it, it does not establish it. That C4 is asserted in prose rather than typed on
the accessor signature is escalation~E4 (\S\ref{sec:conclusion}).

\subsection{CSA margin and collateral}
\label{sec:sec09-csa}

The CSA margin contract is a wallet-level smart contract attached to a per-counterparty collateral
wallet. It reads the aggregate mark-to-market of the trades the CSA governs in the real ledger,
computes the required collateral (threshold, minimum transfer amount, eligible-collateral and
haircut rules), and emits margin-call moves to or from the collateral wallet. Margin eligibility
and haircut are contractual data in \textit{ProductTerms}; the exposure is the observed valuation.

For a synthetic managed account the netting set is the single real unit $u_{\mathrm{TRS}}$, never
the simulated constituents in the virtual ledger: reading those would breach the isolation
invariant (\S\ref{sec:sec09-trs}) and over-margin unmargined positions. The realm tag and C4
realm-scoped reads enforce this. The margin exposure $N_k\cdot\mathrm{TR}_k$ is itself an
observation of the virtual book lifted across the realm boundary as a number; that this
observation is unattested and unversioned is escalation~E2 (\S\ref{sec:conclusion}).

\subsection{TRS as synthetic managed account}
\label{sec:sec09-trs}

A total return swap connects a virtual ledger $\mathcal{L}_v$ to real counterparty wallets in the
real ledger $\mathcal{L}_r$ through periodic cash settlement. Ledger isolation is invariant~P7: no
move has one endpoint in $\mathcal{L}_v$ and the other in $\mathcal{L}_r$. A minimal virtual
ledger is not a second engine but a $\mathrm{realm}\in\{\text{real},\text{virtual}\}$ tag enforcing
P7.

At reset $t_k$ the TRS contract in $\mathcal{L}_r$ observes the virtual book's performance and
emits one real move of $\big|\mathrm{Payment}_k\big|$,
$\mathrm{Payment}_k=N_k\cdot\mathrm{TR}_k-N_k\cdot r_k\cdot\Delta t_k$, direction by sign,
$w_{\text{payer}}\to w_{\text{receiver}}$ in USD, tagged \texttt{TRS\_NET\_SETTLEMENT}. The total
return $\mathrm{TR}_k=(V^v_{t_k}-V^v_{t_{k-1}})/V^v_{t_{k-1}}$ is partial: it requires
$V^v_{t_{k-1}}>0$, a precondition that is a typed failure, not a division by zero.

\begin{theorem}[Equivalence to the funded managed account]
\label{thm:trs-equiv}
The TRS settlement and the funded managed-account fee/PnL settlement are the same ledger
operation: one net conservative cash move whose magnitude is a deterministic function of an
observed performance and a recorded baseline. The funded book plays the role of $\mathcal{L}_v$;
the periodic settlement is the TRS payment.
\end{theorem}
\noindent\textit{Proof sketch.} Both reduce to
$\mathrm{Crystallise}(\text{magnitude},\text{from},\text{to})$ with magnitude $=$ observed
performance over the interval and direction by sign; both reset a baseline to the post-settlement
value; neither moves a position across a boundary. The settlement primitive is one; the products
differ, and that difference is a CDM (reporting) distinction, not a ledger one: a CDM
\texttt{TotalReturnSwap} is two-legged ($N\,\mathrm{TR}$ and $N\,r\,\Delta t$), so the report
derives from the retained CDM \texttt{BusinessEvent}, never from the netted move
(\S\ref{sec:sec09-conformance}). \qed

\paragraph{Price consistency.} $\mathcal{L}_v$ valuation and $\mathcal{L}_r$ settlement use one
price vector $P_t$. Because P7 forbids a cross-ledger move, a price divergence produces unexplained
PnL with no internal reconciliation path --- the one place the no-internal-break guarantee can be
silently voided. Both realms read the price from the same logged event, so a fresh replay
reconstructs one price for both; what remains is attestation of that logged price, escalation~E2,
the Nazarov attestation finding (\S\ref{sec:conclusion}).

\subsection{Redemption and wind-down}
\label{sec:sec09-redemption}

Wind-down strikes a final NAV, settles the final fee, returns capital, and retires the mandate.
Positions are liquidated against the market wallet; the cash is paid to the client's external
wallet; the mandate unit is returned to the manager. All closed-out rows are retained at zero under
the monotone carrier (C1), preserving the audit trail, the final high-water mark, and the
quantities needed for tax reconstruction.
\begin{lstlisting}
Move(from: w_C,     to: w_mkt, unit: <asset>, quantity: <held>)      -- liquidate
Move(from: w_mkt,   to: w_C,   unit: USD,     quantity: <proceeds>)
Move(from: w_C,     to: w_ext, unit: USD,     quantity: <final NAV>) -- capital-out
Move(from: w_C,     to: w_M,   unit: u_MA,    quantity: 1)           -- retire mandate
\end{lstlisting}

\subsection{Account-level balance-sheet substantiation}
\label{sec:sec09-substantiation}

Each account's balance at any date is a deterministic projection of the move stream filtered to its
wallets; client statements, desk PnL, and regulatory position reports differ only in the filter.
Within the ledger there is no separate account-level quantity record to reconcile against. The move
stream substantiates quantities and provenance; valuations depend on the external price vector, and
disclosures and legal status require external evidence. Net settlement is operational netting:
presenting a balance sheet net requires, in addition, a legally enforceable right of set-off and
the intent to settle net (IAS~32.42 / ASC~210-20), a determination outside the ledger. The
accounts' assets~$=$~liabilities~$+$~equity reconciliation holds at the quantity level by P1 (the
trial balance balances by construction); the valued and classified presentation does not follow
from P1 alone.

\subsection{Worked example}
\label{sec:sec09-example}

One managed account $w_C$ under mandate $u_{\mathrm{MA}}$ issued by manager $w_M$; virtual market
wallet $w_{\text{mkt}}$; client external wallet $w_{\text{ext}}$. USD reference. Terms: management
fee $2\%/\text{yr}$, performance fee $20\%$ over a high-water mark (no hurdle), quarterly
($\Delta t=0.25$). Conventions: management fee on end-of-period NAV; performance fee on NAV net of
the management fee, above the mark; the mandate is non-valued. $\sum_w w(u)=0$ holds for every unit
after every move.

\begin{longtable}{p{4.3cm}p{7.2cm}r}
\toprule
\textbf{Step} & \textbf{Moves} ($\text{src}\to\text{dst}:\text{qty unit}$) & \textbf{NAV} \\
\midrule
A. Mandate issuance & $w_M\to w_C:1\ u_{\mathrm{MA}}$ & --- \\
B. Subscription & $w_{\text{ext}}\to w_C:1{,}000{,}000\ \text{USD}$ & $1{,}000{,}000$ \\
C. Trade (buy $5{,}000$ AAPL @100) & $w_C\to w_{\text{mkt}}:500{,}000\ \text{USD}$;\ \ $w_{\text{mkt}}\to w_C:5{,}000\ \text{AAPL}$ & $1{,}000{,}000$ \\
D. Q1 NAV (AAPL $\to 130$) & (price move; no ledger move) & $\mathbf{1{,}150{,}000}$ \\
E. Fee crystallisation & $w_C\to w_M:5{,}750\ \text{USD}$ (mgmt);\ \ $w_C\to w_M:28{,}850\ \text{USD}$ (perf) & $\mathbf{1{,}115{,}400}$ \\
F. Wind-down & $w_C\to w_{\text{mkt}}:5{,}000\ \text{AAPL}$;\ \ $w_{\text{mkt}}\to w_C:650{,}000\ \text{USD}$;\ \ $w_C\to w_{\text{ext}}:1{,}115{,}400\ \text{USD}$;\ \ $w_C\to w_M:1\ u_{\mathrm{MA}}$ & $0$ \\
\bottomrule
\end{longtable}

\paragraph{Fee arithmetic (Step~E).} No capital flow occurred inside the period, so
$\mathrm{Perf}=(1{,}150{,}000-1{,}000{,}000)-0=150{,}000$. Management fee
$=1{,}150{,}000\times0.02\times0.25=5{,}750$. NAV net of management $=1{,}144{,}250$. Performance
fee $=0.20\times(1{,}144{,}250-1{,}000{,}000)=28{,}850$. NAV after both fees $=1{,}115{,}400$. New
high-water mark $=\max(1{,}000{,}000,\,1{,}115{,}400)=1{,}115{,}400$; reset baseline
$B_1=1{,}115{,}400$. Fee zero-sum (Theorem~\ref{thm:fee-zerosum}):
$\sum_w\Delta_{\text{fee}}(\text{USD})=-34{,}600\,(w_C)+34{,}600\,(w_M)=0$.

\paragraph{Closing identity (all wallets).} Final non-zero balances: $w_M=+34{,}600$ (fee revenue),
$w_{\text{mkt}}=-150{,}000$ (the gross gain paid out by the market), $w_{\text{ext}}=+115{,}400$
(client net). Their sum is $34{,}600-150{,}000+115{,}400=0$. The client received principal
$1{,}000{,}000$ plus net gain $115{,}400$; the gross gain $150{,}000$ equals net $115{,}400$ plus
fees $34{,}600$. Per-unit conservation holds across the chain: $u_{\mathrm{MA}}$ is issued
$(+1,-1)$ then returned $(0,0)$; USD nets to $0$ at $w_C$ and globally after every move; AAPL is
acquired and liquidated with $w_{\text{mkt}}$ as the mirror. Every position row is retained at zero
under the monotone carrier.

\subsection{Conformance and flags}
\label{sec:sec09-conformance}

No conformance is asserted; each item is demonstrated against a cited source or flagged for the
responsible owner. The CDM strata named below (\texttt{LegalAgreement}, \texttt{Trade},
\texttt{TotalReturnSwap}, \texttt{TradeState}, \texttt{BusinessEvent}) are identified at the type
level only; pinning them to a named CDM release is owed to the CDM owner and is flagged, not
asserted.

\begin{itemize}[leftmargin=*]
\item \textbf{CDM mandate representation.} $u_{\mathrm{MA}}$ maps to the CDM
\texttt{LegalAgreement} stratum (the ISDA Create-digitised investment-management agreement:
parties, fee schedule, benchmark identity, methodology, limits), not the \texttt{Trade}/product
stratum. A discretionary mandate is a legal agreement, not a tradable contract.
\item \textbf{CDM TRS representation.} A CDM \texttt{TotalReturnSwap} is two-legged; the report is
the retained \texttt{BusinessEvent} (performance leg and financing leg $N\,r\,\Delta t$), never the
netted settlement move. The ledger ``same mechanism'' identity (Theorem~\ref{thm:trs-equiv}) is a
settlement-primitive identity, flagged as not a CDM/reporting identity.
\item \textbf{SFTR/EMIR reporting surface.} $u_{\mathrm{MA}}$ issuance is not reportable: a
discretionary mandate is the MiFID~II Annex~I~\S~A(4) portfolio-management service, not a Section~C
instrument, not an SFT, not a derivative.
The reportable artefact is never $u_{\mathrm{MA}}$ but the underlying book trade or
$u_{\mathrm{TRS}}$, each carrying its own UTI/UPI/LEI. Reportability is a per-relationship,
regime-cited attribute on the wrapped unit's \textit{ProductTerms}, not inferable from the coarse
unit-type key. Governance of that attribute, together with the dual-sided EMIR/CFTC/SFTR reporting
that needs validated ISO-17442 counterparty LEIs and classification---neither quantity facts nor
with any home in the move stream---is escalation~E3 (\S\ref{sec:conclusion}).
\item \textbf{IFRS classification.} Economic NAV and economic PnL are derived here; accounting
classification of fair-value changes (FVTPL/FVOCI/amortised cost) and net presentation
(IAS~32.42) are outside scope and require external evidence. Accrued fees, modelled as conserved
moves rather than stored scalars, carry their own contra-entry and are auditable as double-entry.
\item \textbf{CDM \texttt{TradeState} alignment.} That the CDM
\texttt{TradeState}-per-\texttt{Trade} view and \textit{PositionState}$[w,u]$ stay replay-consistent
(both projections of one minted event) is asserted, not yet verified; it is the outstanding F6
deliverable (a Rosetta mapping rerun against the three-map schema).
\end{itemize}

Escalations E1--E5 and the Nazarov attestation finding are cited at their points of origin above
and consolidated in the flagged register (\S\ref{sec:conclusion}). They are not failures of the workflow,
which is sound as quantity algebra --- every step is atomic $q>0$ moves with $\sum_w w(u)=0$
verified, and the worked example closes to the penny --- but the points at which the performance and
reporting layer outruns the framework's current treatment of recorded observations and external
identity.
